_{\mathrm{sample}}(z)
    }{
        \sum_{z\in\mathcal N}
        I_{\mathrm{sample}}(z)
    }.
    \label{eq:near_retention}
\end{equation}
Taper parameters are calibrated on training contexts so that the
compared selective rules retain the same target near-field fraction,
up to the predetermined tolerance. Parameters are then frozen before
evaluation. This comparison isolates how the rules treat the far-field
tail after preserving comparable local utility.

\paragraph{Aggregate-budget matching.}
For a weighting function $\omega$, define its detached negative-update
budget as
\begin{equation}
    B(\omega)
    =
    \sum_{z\in\mathcal B^-}
    \omega(D_\theta(z))
    I_{\mathrm{sample}}(z).
    \label{eq:taper_budget}
\end{equation}
The budget-matched global coefficient is
\begin{equation}
    \alpha_{\mathrm{budget}}
    =
    \frac{
        B(\omega)
    }{
        B(\omega_{\mathrm{uncontrolled}})
    }.
    \label{eq:taper_global_match}
\end{equation}
This control applies the same total first-order negative-update budget
without using remoteness. A selective taper can therefore be credited
only for improvements that remain after comparison with its
budget-matched global baseline.

\paragraph{Measurements.}
Each method is evaluated over the full prespecified training horizon.
We report task performance, held-out-context generalization, distance
from the hidden optimum, policy displacement, entropy or variance,
effective support, near-field retention, aggregate negative-update
budget, and all three terminal-event categories. Method comparisons use
the same seeds and paired statistics.



\subsection{External Taper Validation}
\label{app:external_taper_protocol}

\paragraph{Scientific scope.}
External taper experiments test whether the controlled weighting
principle transfers to complex policies. They do not reproduce the exact
counterfactual geometry of C-U1 or D-U1 and therefore do not replace the
controlled source and causal experiments.

\paragraph{Hopper.}
All Hopper taper methods branch from the same actor initialization and
use the same offline transitions and advantage labels. For an action
with inverse-squash coordinate $u$, Gaussian mean $\mu_\theta(s)$, and
diagonal scale $\sigma_\theta(s)$, the remoteness statistic is
\begin{equation}
    D_{\mathrm{Hopper}}(s,a)
    =
    \frac{1}{2}
    \sum_j
    \left(
        \frac{
            u_j-\mu_{\theta,j}(s)
        }{
            \sigma_{\theta,j}(s)
        }
    \right)^2.
    \label{eq:hopper_taper_remoteness}
\end{equation}
The additive Gaussian normalization term is omitted because it does not
change the ordering used by the taper. Remoteness is recomputed under
the current actor and detached before constructing the negative weight.
The mechanism protocol with frozen advantages is kept separate from the
standard offline-RL performance protocol; the former studies transfer of
the far-field control mechanism, whereas the latter reports final
environment return.

\paragraph{Countdown.}
The Countdown offline dataset is generated once from a frozen SFT
policy. For each puzzle, it contains a verified positive response, a
near negative response with low mean token surprisal, and a far negative
response with high mean token surprisal. Mean token surprisal is
\begin{equation}
    D_{\mathrm{Countdown}}(x,y)
    =
    -\frac{1}{|y|}
    \sum_{t=1}^{|y|}
    \log
    \pi_\theta
    \left(
        y_t
        \mid
        x,y_{<t}
    \right).
    \label{eq:countdown_taper_remoteness}
\end{equation}
During method training, the responses remain fixed but their surprisal is
recomputed under the current policy. The sequence-level taper weight is
computed from the detached current surprisal and multiplies the negative
sequence objective.

All Countdown methods start from the same SFT checkpoint and use the
same positive, near-negative, and far-negative responses, mixture
coefficients, validation set, and maximum training budget. Checkpoint
selection uses the registered validation task metric, and final verifier
success is measured only on the held-out test puzzles.

\paragraph{Comparisons and reporting.}
The common external comparisons are Positive-only, Uncontrolled,
Global scaling, and exponential DRPO. Domain-specific baselines are
reported in the overall-performance section rather than treated as
controlled taper variants. Hopper reports normalized return and policy
statistics; Countdown reports greedy verifier success, pass@$k$, valid
expression rate, surprisal, and support statistics. Task-performance,
boundary, and NaN/Inf outcomes remain separate in both domains.
\subsection{Shared Objectives and Negative-Update Controls}
\label{app:shared_methods}

This section defines the protocol controls and remoteness-aware
weighting rules shared by the D4RL and Countdown experiments. Their
task-specific remoteness signals, data construction, and optimization
protocols are described separately in
Appendices~\ref{app:d4rl_methods} and
\ref{app:countdown_methods}.

\paragraph{Shared signed objective.}

Let \(z\) denote an offline transition or response, let
\(A(z)\) denote its fixed signed learning signal during an actor update,
and let \(D_\theta(z)\) denote its learner-relative remoteness under the
current policy. We write

\[
A^{+}(z)=\max\{A(z),0\},
\qquad
A^{-}(z)=\max\{-A(z),0\}.
\]

The shared family of actor objectives is

\[
\mathcal{J}_{\omega}(\theta)
=
\mathbb{E}_{z\sim\mathcal{B}}
\left[
A^{+}(z)\log \pi_\theta(z)
 -
\omega\!\left(D_\theta(z)\right)
A^{-}(z)\log \pi_\theta(z)
\right],
\]

where positive updates retain unit weight and
\(\omega(D)\) controls only the negative branch. During each actor
update, the advantage and remoteness-derived weight are treated as
fixed coefficients; gradients are not propagated through the weighting
rule.

For the remoteness-aware methods, define the normalized excess
remoteness

\[
x_\theta(z)
=
\left[
\frac{D_\theta(z)-\tau}{c}
\right]_{+},
\]

where \(\tau\) is the near-field threshold, \(c>0\) is the task-specific
scale, and \([u]_{+}=\max\{u,0\}\). Consequently, samples inside the
near-field region \(D_\theta(z)\leq\tau\) retain their full negative
weight.

\paragraph{Positive-only.}

Positive-only removes all negative updates:

\[
\omega_{\mathrm{pos}}(D)=0.
\]

It therefore provides an extreme no-repulsion reference. In D4RL it
retains only transitions with positive actor advantages; in Countdown
it retains only verified successful responses under the corresponding
binary learning signal.

\paragraph{Global-\(\alpha\).}

Global-\(\alpha\) applies the same attenuation to every negative sample:

\[
\omega_{\mathrm{global}}(D)=\alpha,
\qquad
0\leq\alpha\leq 1.
\]

This control tests whether performance improvements can be explained by
a global reduction in negative-update magnitude, without using
learner-relative remoteness.

\paragraph{Reciprocal-Linear and Reciprocal-Quadratic.}
For the remoteness-aware methods, define the normalized excess
remoteness
\[
x_\theta(z)
=
\left[
\frac{D_\theta(z)-\tau}{c}
\right]_{+},
\]
and its radialized coordinate
\[
r_\theta(z)=\sqrt{x_\theta(z)}.
\]
For Gaussian policies, \(r_\theta\) is asymptotically proportional to
standardized distance in the far field; for sequence policies, it is a
surprisal-derived radial coordinate rather than a geometric distance in
the discrete action space.
The two reciprocal controls use the same remoteness signal and
near-field threshold as DRPO, but differ in their far-field decay:

\[
\omega_{\mathrm{Rec\text{-}L}}(D)
=
\frac{1}{1+\lambda r_\theta(z)},
\]

and

\[
\omega_{\mathrm{Rec\text{-}Q}}(D)
=
\frac{1}{1+\lambda r_\theta(z)^2}
=
\frac{1}{1+\lambda x_\theta(z)}.
\]

Reciprocal-Linear has the slowest far-field attenuation. The quadratic
variant decays more rapidly but retains a polynomial tail. These
controls isolate whether the benefit comes merely from using a
remoteness signal or specifically from the decay profile selected by
DRPO.

\paragraph{DRPO.}

DRPO uses an exponential remoteness taper:

\[
\omega_{\mathrm{DRPO}}(D)
=
\exp\left\{
-\lambda x_\theta(z)
\right\}.
\]

The weight equals one throughout the near-field region and then decays
exponentially once remoteness exceeds \(\tau\). Thus, DRPO preserves
local negative feedback while suppressing the remote negative tail more
rapidly than the reciprocal alternatives.

\paragraph{Hyperparameter selection.}

Within each task, \(\tau\), \(c\), \(\lambda\), and the
Global-\(\alpha\) coefficient are selected using the registered
validation protocol only. Test-task performance is not used for
hyperparameter selection. All protocol-matched methods share the same
candidate budgets and checkpoint-selection rule. The task-specific
search spaces and selected values are reported in
Appendices~\ref{app:d4rl_methods} and
\ref{app:countdown_methods}.


\subsection{D4RL Baselines and Experimental Protocol}
\label{app:d4rl_methods}

The definitions of Positive-only, Global-\(\alpha\),
Reciprocal-Linear, Reciprocal-Quadratic, and DRPO follow
Appendix~\ref{app:shared_methods}. This section specifies only the
D4RL-specific baselines, remoteness construction, training protocol,
and evaluation procedure.

\paragraph{Tasks and datasets.}

We evaluate HalfCheetah, Hopper, and Walker2d on the medium,
medium-replay, and medium-expert datasets, producing the nine tasks
reported in Table~\ref{tab:d4rl_performance}. The same fixed offline
dataset is used by every protocol-matched method within a task.

\paragraph{Published offline-RL baselines.}

CQL, TD3+BC, and IQL are included as published reference results.
These values are not protocol-matched reruns and are therefore marked
by \(\dagger\) in Table~\ref{tab:d4rl_performance}. The main table
reports their published task-level means. Their original uncertainty
statistics and exact score provenance are recorded separately from the
protocol-matched DRPO experiments.

\paragraph{Gaussian-policy remoteness.}

For an action \(a\in\mathbb{R}^{d_a}\) and current Gaussian policy with
mean \(\mu_\theta(s)\) and standard deviation \(\sigma_\theta(s)\), we use the
following implementation-level standardized remoteness statistic

\[
D_{\mathrm{D4RL}}(s,a)
=
\frac{1}{d_a}
\sum_{j=1}^{d_a}
\left(
\frac{a_j-\mu_{\theta,j}(s)}
     {\sigma_{\theta,j}(s)}
\right)^2.
\]

This dimension-normalized convention differs from
Equation~\eqref{eq:hopper_taper_remoteness} only by a fixed task-level scale.
The associated threshold and scale parameters absorb that constant, and each
protocol keeps its convention fixed across all compared methods.

When the actor uses a transformed action distribution, this quantity is
computed in the policy's native Gaussian coordinates. The resulting
remoteness is used only to determine the negative-update coefficient;
the distance term is detached during the actor update.

\paragraph{Protocol-matched actor comparison.}

Positive-only, Global-\(\alpha\), Reciprocal-Linear,
Reciprocal-Quadratic, and DRPO share the same offline dataset, critic
checkpoint or critic-training protocol, actor initialization, network
architecture, optimizer, minibatch schedule, actor-update budget,
evaluation schedule, and random seeds. They differ only in the
negative-update weighting rule defined in
Appendix~\ref{app:shared_methods}.

The common actor advantage construction is held fixed across these
methods. No method is allowed to change the critic, relabel samples, or
alter the data distribution in order to improve its actor result.

\paragraph{Hyperparameters.}

The registered training and selection configuration is summarized in
Table~\ref{tab:d4rl_registered_config}. All entries must be populated
directly from the frozen experiment configuration rather than
reconstructed after observing the test results.

\begin{table}[t]
\centering
\caption{Registered configuration for the D4RL performance experiments.}
\label{tab:d4rl_registered_config}

\small
\setlength{\tabcolsep}{8pt}
\renewcommand{\arraystretch}{1.12}

\begin{tabular}{@{}ll@{}}
\toprule
Configuration & Registered value \\
\midrule
Actor-update budget
& \textbf{[registered value]} \\
Evaluation interval
& \textbf{[registered value]} \\
Evaluation episodes
& \textbf{[registered value]} \\
Random seeds
& \textbf{[registered seed set]} \\
Near-field threshold candidates, $\tau$
& \textbf{[registered grid]} \\
Remoteness-scale candidates, $c$
& \textbf{[registered grid]} \\
Taper-strength candidates, $\lambda$
& \textbf{[registered grid]} \\
Global-scaling candidates, $\alpha$
& \textbf{[registered grid]} \\
\bottomrule
\end{tabular}
\end{table}

\paragraph{Evaluation and uncertainty.}

Task performance is reported using the standard D4RL normalized return.
For every protocol-matched method, the main table reports the mean over
the registered seeds, while seed-level values and uncertainty
statistics are reported in the appendix. The D4RL-9 total is the sum of
the nine task-level means.

Task-performance collapse is assessed from normalized return and
training trajectories. Variance or support-boundary events are reported
separately from task performance. NaN/Inf numerical failures are
reported as a third outcome and are not merged with either low return or
policy-boundary events. Any claim concerning convergence, collapse, or
method ranking is based on the registered terminal audit rather than an
intermediate checkpoint.


\subsection{Countdown Baselines and Experimental Protocol}
\label{app:countdown_methods}

Positive-only, Global-\(\alpha\), Reciprocal-Linear,
Reciprocal-Quadratic, and DRPO follow the shared definitions in
Appendix~\ref{app:shared_methods}. This section defines the two
published off-policy baselines and the Countdown-specific response-bank,
remoteness, training, and evaluation protocols.

\paragraph{AsymRE.}

Asymmetric REINFORCE defines the response-level advantage using a
tunable reward baseline \(V\):

\[
A_{\mathrm{AsymRE}}(x,y)
=
r(x,y)-V,
\]

and optimizes

\[
\mathcal{J}_{\mathrm{AsymRE}}(\theta)
=
\mathbb{E}_{(x,y)\sim\mathcal{B}}
\left[
\bigl(r(x,y)-V\bigr)
\log\pi_\theta(y\mid x)
\right].
\]

For a binary verifier reward, lowering \(V\) emphasizes successful
responses, whereas increasing \(V\) strengthens suppression of failed
responses. The baseline is selected using the validation protocol and
is held fixed within each run.

Global-\(\alpha\) and AsymRE remain separate comparisons because they
operate on different underlying learning signals. Global-\(\alpha\)
scales the negative part of the common precomputed signed advantage,
whereas AsymRE constructs its signed advantage directly from the raw
verifier reward and the published baseline parameterization. Under
binary rewards and simplified normalization they are closely related,
but the formal comparison preserves their native objectives and tuning
rules.

\paragraph{TOPR.}

Let \(\mu(y\mid x)\) denote the frozen response-generating policy and

\[
\rho_\theta(x,y)
=
\frac{\pi_\theta(y\mid x)}
     {\mu(y\mid x)}
\]

the sequence-level i